\documentclass[11pt,a4paper]{article}

\usepackage[margin=1in]{geometry}
\usepackage{amsmath,amssymb,amsthm}
\usepackage{mathtools}
\usepackage{enumerate}
\usepackage{hyperref}
\usepackage{cleveref}
\usepackage[ruled,vlined,linesnumbered]{algorithm2e}

\newtheorem{theorem}{Theorem}[section]
\newtheorem{lemma}[theorem]{Lemma}
\newtheorem{proposition}[theorem]{Proposition}
\newtheorem{corollary}[theorem]{Corollary}
\newtheorem{definition}[theorem]{Definition}
\newtheorem{remark}[theorem]{Remark}

\DeclareMathOperator{\proj}{proj}
\DeclareMathOperator{\AM}{AM}
\DeclareMathOperator{\HM}{HM}
\DeclareMathOperator{\diag}{diag}
\DeclareMathOperator{\tr}{tr}
\DeclareMathOperator{\rank}{rank}
\newcommand{\R}{\mathbb{R}}
\newcommand{\E}{\mathbb{E}}
\newcommand{\Prob}{\mathbb{P}}
\newcommand{\ip}[2]{\langle #1, #2 \rangle}
\newcommand{\norm}[1]{\|#1\|}

\title{The Blessing of Dimensionality for Policy Gradients:\\
Sparse Recovery, Evidence Weighting, and the\\
Logarithmic Cost of Action Space Expansion}
\author{Eugene Shcherbinin}
\date{March 2026}

\begin{document}
\maketitle

\begin{abstract}
The REINFORCE gradient estimator suffers from variance that appears to
scale linearly with the action space dimension $d = |\mathcal{A}|$. We show
this is a curse of dimensionality illusion: when the optimal policy is
$k$-sparse (concentrates on $k \ll d$ actions), high-dimensional concentration
phenomena reduce the effective sample complexity to $O(k \log(d/k))$ ---
\emph{sub-linear} in~$d$. We formalize this through three mechanisms:
(1)~measure concentration bounds showing gradient noise self-averages
across irrelevant action dimensions; (2)~compressed sensing theory giving
sparse gradient recovery guarantees; (3)~an effective-dimension refinement
of the evidence weight from the $\Omega$-framework that captures the true
information content of the agent's experience. Combined with L1 policy
regularization, this yields a natural curriculum: start sparse (low
variance, fast convergence), expand as evidence improves (logarithmic
cost). The blessing of dimensionality transforms the action space expansion
problem from multiplicative to logarithmic.
\end{abstract}


\section{The Variance Problem in Large Action Spaces}

\subsection{REINFORCE and the Apparent Curse}

The REINFORCE gradient estimator for a policy $\pi_\theta$ over
action space $\mathcal{A}$ with $|\mathcal{A}| = d$ is:
\begin{equation}\label{eq:reinforce}
  \hat{g}(\tau) = R(\tau) \cdot \nabla_\theta \log \pi_\theta(\tau)
  \in \R^d.
\end{equation}
The variance of this estimator decomposes by action:
\begin{equation}\label{eq:variance-decomp}
  \text{Var}(\hat{g}) = \sum_{a=1}^{d}
    \underbrace{\pi(a)}_{\text{sampling prob.}} \cdot
    \underbrace{\left(\frac{\nabla_\theta \pi(a)}{\pi(a)}\right)^2}_{\text{score}^2} \cdot
    \underbrace{R(a)^2}_{\text{reward}^2}
    \;-\; \norm{\nabla J}^2.
\end{equation}

For a low-probability action $a$ with $\pi(a) = \epsilon \to 0$:
\begin{equation}
  \text{contribution}_a \;\propto\; \frac{R(a)^2}{\pi(a)} \;\to\; \infty.
\end{equation}

This is the \textbf{inverse probability weighting} problem: rare actions
contribute enormous variance spikes when sampled. Expanding $\mathcal{A}$
by adding many low-probability actions appears to make this sum diverge.

\textbf{The apparent conclusion}: variance scales as $\Theta(d)$, so
REINFORCE is impractical for large action spaces.

\textbf{We show this conclusion is wrong} when the optimal policy is sparse.


\subsection{The Overfitting Analogy}

The parallel to classical ML overfitting is precise:

\begin{center}
\begin{tabular}{lll}
\hline
& \textbf{Supervised learning} & \textbf{Policy gradient} \\
\hline
Space     & Parameter space $\R^d$ & Action space $\mathcal{A}$, $|\mathcal{A}| = d$ \\
Problem   & Too many params vs.\ data & Too many actions vs.\ samples \\
Noise     & Fitting noise in labels & Gradient spikes from rare actions \\
Signal    & True function (low-dim) & True gradient (sparse) \\
Cure      & L1/Lasso (sparsity) & L1 on policy (action sparsity) \\
Blessing  & Compressed sensing & Sparse gradient recovery \\
\hline
\end{tabular}
\end{center}

The key insight from compressed sensing: \emph{sparsity makes
high-dimensional problems easier, not harder.}


\section{The Blessing: Concentration and Sparse Recovery}

\subsection{Concentration of the Gradient Vector}

\begin{theorem}[Gradient concentration in high dimensions]\label{thm:concentration}
Let $\hat{g}_1, \ldots, \hat{g}_M$ be i.i.d.\ REINFORCE gradient samples
in $\R^d$. Suppose the gradient covariance $\Sigma = \text{Cov}(\hat{g})$
has eigenvalues $\lambda_1 \geq \cdots \geq \lambda_d$. Then:
\begin{equation}\label{eq:concentration}
  \Prob\!\left(\left\|\frac{1}{M}\sum_{m=1}^M \hat{g}_m - \nabla J\right\|
  \leq \epsilon\right) \geq 1 - \delta
\end{equation}
provided $M \geq C \cdot \frac{\tr(\Sigma)}{\epsilon^2} \cdot \log(d/\delta)$,
where $C$ is a universal constant.
\end{theorem}

\begin{proof}[Proof sketch]
Apply the matrix Bernstein inequality to the sum of centered random
vectors $\hat{g}_m - \nabla J$. The key quantity is $\tr(\Sigma)$
(not $d \cdot \lambda_{\max}$) because the noise contributions from
different eigenspaces are independent and add their variances, not
their maxima.
\end{proof}

\begin{remark}[Why $\tr(\Sigma)$, not $d \cdot \sigma^2$]
The naive bound treats all $d$ dimensions equally:
$M \geq d \cdot \sigma^2_{\max} / \epsilon^2$.
But $\tr(\Sigma) = \sum_i \lambda_i$ can be \emph{much} smaller than
$d \cdot \lambda_1$ when the eigenspectrum decays. This is the first
blessing: gradient variance concentrates on a low-dimensional subspace.
\end{remark}


\subsection{Effective Dimension}

\begin{definition}[Effective dimension of the gradient]\label{def:deff}
The effective dimension of the policy gradient at $\pi_\theta$ is:
\begin{equation}\label{eq:deff}
  d_{\text{eff}} \;=\; \frac{\tr(\Sigma)}{\lambda_{\max}(\Sigma)}
  \;=\; \frac{\sum_{i=1}^d \lambda_i}{\lambda_1}.
\end{equation}
\end{definition}

This measures how many eigenvalues contribute meaningfully.
If only $k$ eigenvalues are large and the rest are negligible:
$d_{\text{eff}} \approx k$. The sample complexity
(\cref{thm:concentration}) becomes:
\begin{equation}
  M \;\geq\; C \cdot \frac{d_{\text{eff}} \cdot \lambda_1}{\epsilon^2}
  \cdot \log(d/\delta)
  \;\approx\; O\!\left(\frac{k}{\epsilon^2} \log d\right).
\end{equation}

\textbf{This is sub-linear in $d$.} Adding 10,000 low-probability actions
to a 10-action game costs an extra $\log(10000/10) \approx 7$ factor,
not a $1000\times$ factor.


\subsection{Sparse Policy Recovery via Compressed Sensing}

When the optimal policy $\pi^*$ is $k$-sparse (at most $k$ actions
have positive probability in each state), compressed sensing theory
gives sharp guarantees.

\begin{theorem}[Sparse gradient recovery]\label{thm:cs}
Let $\pi^*$ be $k$-sparse. Under the restricted isometry property
(RIP) on the gradient measurement matrix, the true gradient direction
can be recovered from $M$ REINFORCE samples with:
\begin{equation}\label{eq:cs-bound}
  \norm{\hat{\nabla} J - \nabla J} \leq C_1 \cdot \frac{\sigma_k(\nabla J)_1}{\sqrt{k}}
  + C_2 \cdot \frac{\epsilon}{\sqrt{M}}
\end{equation}
where $\sigma_k(\nabla J)_1$ is the $\ell_1$ error of the best
$k$-sparse approximation, provided:
\begin{equation}
  M \;\geq\; C \cdot k \log(d/k).
\end{equation}
\end{theorem}

\begin{proof}[Proof sketch]
The REINFORCE samples $\hat{g}_m$ are random linear measurements of
$\nabla J$ corrupted by noise. When $\pi^*$ is $k$-sparse, $\nabla J$
is approximately $k$-sparse (gradient is zero in directions corresponding
to zero-probability actions). Apply the Cand\`es--Tao RIP theory to the
measurement ensemble.
\end{proof}

\begin{remark}[The $k \log(d/k)$ scaling]
This is the compressed sensing miracle: you need samples proportional
to the sparsity $k$, not the ambient dimension $d$. The $\log(d/k)$
factor is the information-theoretic price of not knowing \emph{which}
$k$ actions are relevant.
\end{remark}


\subsection{Johnson--Lindenstrauss for Gradient Updates}

\begin{proposition}[Random projection preserves gradient structure]\label{prop:jl}
Let $\Phi \in \R^{m \times d}$ be a random Gaussian matrix with
$m = O(k \log d / \epsilon^2)$. Then with high probability, for all
$k$-sparse gradient vectors $g, g'$:
\begin{equation}
  (1 - \epsilon)\norm{g - g'} \leq \norm{\Phi g - \Phi g'} \leq (1 + \epsilon)\norm{g - g'}.
\end{equation}
\end{proposition}

\textbf{Implementation}: Project $\hat{g}_m \in \R^d$ to
$\tilde{g}_m = \Phi \hat{g}_m \in \R^m$, average in the low-dimensional
space, then recover via sparse recovery (L1 minimization / LASSO).
This reduces the per-step computational cost from $O(d)$ to $O(m) = O(k \log d)$.


\section{Evidence Weighting with Effective Dimension}

\subsection{Refined Evidence Weight}

The EW-PG framework (companion paper) uses the evidence weight:
\begin{equation}
  w_i = \frac{V_{\min}}{V_i}, \qquad V_i = \text{Var}(\hat{g}_i).
\end{equation}

This treats the variance as a scalar, summing over all $d$ dimensions.
But the blessing of dimensionality reveals that most of this variance
is irrelevant noise in unused action dimensions.

\begin{definition}[Effective evidence weight]\label{def:eff-evidence}
The effective evidence weight refines $V_i$ using the gradient's
effective dimension:
\begin{equation}\label{eq:eff-weight}
  V_i^{\text{eff}} = \frac{d_{\text{eff},i}}{d} \cdot V_i,
  \qquad
  w_i^{\text{eff}} = \frac{V_{\min}^{\text{eff}}}{V_i^{\text{eff}}}.
\end{equation}
\end{definition}

\begin{theorem}[Effective evidence improvement]\label{thm:eff-evidence}
For an agent with $k$-sparse optimal policy in a $d$-action game:
\begin{equation}
  \frac{w_i^{\text{eff}}}{w_i} \;=\; \frac{d}{d_{\text{eff},i}} \;\geq\; \frac{d}{k}.
\end{equation}
The effective evidence weight is a factor of $d/k$ larger than the
ambient weight. For $d = 10{,}000$, $k = 5$: the effective weight
is $2{,}000\times$ larger.
\end{theorem}

\begin{remark}[What this means]
An agent in a 10,000-action game with a 5-sparse policy has been
\emph{systematically underweighted} by the ambient evidence measure.
Most of its ``variance'' is noise in action dimensions it will never use.
The effective weight corrects this, allowing the agent to take larger
steps and converge faster.
\end{remark}


\subsection{Estimating Effective Dimension}

\begin{algorithm}[H]
\caption{Effective Dimension Estimation}\label{alg:deff}
\KwIn{Gradient samples $\hat{g}_1, \ldots, \hat{g}_M \in \R^d$}
\KwIn{Threshold $\tau \in (0, 1)$ (e.g., $\tau = 0.95$)}
\KwOut{Effective dimension $\hat{d}_{\text{eff}}$}
\BlankLine
Compute sample covariance $\hat{\Sigma} = \frac{1}{M}\sum_m (\hat{g}_m - \bar{g})(\hat{g}_m - \bar{g})^\top$\;
\BlankLine
\tcp{Option A: Ratio estimator (cheap, $O(Md)$)}
$\hat{d}_{\text{eff}} \leftarrow \tr(\hat{\Sigma}) / \lambda_{\max}(\hat{\Sigma})$\;
\tcp{$\lambda_{\max}$ via power iteration, $O(Md)$}
\BlankLine
\tcp{Option B: Eigenvalue counting (precise, $O(d^2 M + d^3)$)}
Compute eigenvalues $\hat{\lambda}_1 \geq \cdots \geq \hat{\lambda}_d$\;
$\hat{d}_{\text{eff}} \leftarrow \min\{k : \sum_{i=1}^k \hat{\lambda}_i \geq \tau \cdot \tr(\hat{\Sigma})\}$\;
\BlankLine
\tcp{Option C: Random projection (scalable, $O(Mm)$)}
Project: $\tilde{g}_m = \Phi \hat{g}_m$ with $\Phi \in \R^{m \times d}$, $m \ll d$\;
$\hat{d}_{\text{eff}} \leftarrow \tr(\tilde{\Sigma}) / \lambda_{\max}(\tilde{\Sigma})$ in the projected space\;
\end{algorithm}


\section{L1-Regularized Evidence-Weighted PG}

\subsection{The Sparse EW-PG Update}

Combining L1 regularization (for sparsity) with evidence weighting
(for adaptive step sizes) and the effective dimension (for accurate evidence):

\begin{algorithm}[H]
\caption{Sparse Evidence-Weighted PG (Sparse-EW-PG)}\label{alg:sparse-ewpg}
\KwIn{Game $\mathcal{G}$, $N$ agents, $|\mathcal{A}_i| = d_i$}
\KwIn{Learning rates $\gamma_n$, L1 coefficient $\mu > 0$}
\KwIn{Sparsity relaxation schedule $\mu_n = \mu / (1 + n/n_0)$}
\KwOut{Sparse policies converging to Nash}
\BlankLine
Initialize $\pi_i \in \Pi$ for each agent\;
\BlankLine
\For{$n = 1, 2, \ldots$}{
  \BlankLine
  \tcp{--- Gradient and evidence ---}
  \For{$i = 1, \ldots, N$}{
    Compute $\hat{v}_{i,n}$ (REINFORCE)\;
    Update gradient covariance estimate $\hat{\Sigma}_{i,n}$\;
    Estimate effective dimension $d_{\text{eff},i,n}$ (Algorithm~\ref{alg:deff})\;
    \BlankLine
    Compute effective evidence weight:\;
    $V_{i,n}^{\text{eff}} \leftarrow (d_{\text{eff},i,n} / d_i) \cdot \hat{V}_{i,n}$\;
    $w_{i,n}^{\text{eff}} \leftarrow V_{\min,n}^{\text{eff}} / V_{i,n}^{\text{eff}}$\;
  }
  \BlankLine
  \tcp{--- Sparse update with proximal step ---}
  \For{$i = 1, \ldots, N$}{
    Gradient step:\;
    $\tilde{\pi}_{i} \leftarrow \pi_{i,n} + \gamma_n \cdot w_{i,n}^{\text{eff}} \cdot \hat{v}_{i,n}$\;
    \BlankLine
    L1 proximal step (soft-thresholding):\;
    \For{each action $a$}{
      $\tilde{\pi}_i(a) \leftarrow \text{sign}(\tilde{\pi}_i(a)) \cdot \max(|\tilde{\pi}_i(a)| - \gamma_n \mu_n, \; 0)$\;
    }
    \BlankLine
    Project back to simplex:\;
    $\pi_{i,n+1} \leftarrow \proj_{\Delta}(\tilde{\pi}_i)$\;
  }
  \BlankLine
  \tcp{--- Adaptive sparsity relaxation ---}
  \tcp{As evidence improves, relax L1 to allow exploration}
  \tcp{$\mu_n$ decreases: start sparse, expand action support}
}
\end{algorithm}

\subsection{The Natural Curriculum}

The interaction between L1 regularization and evidence weighting creates
a \textbf{self-organizing curriculum}:

\begin{enumerate}
  \item \textbf{Early training} ($n$ small, $\mu_n$ large, $V_i$ high):
    \begin{itemize}
      \item L1 drives policy toward extreme sparsity ($k \approx 1$--$2$)
      \item $d_{\text{eff}} \approx k \ll d$ → effective evidence is high
      \item High $w_i^{\text{eff}}$ → large steps → fast convergence to sparse optimum
    \end{itemize}

  \item \textbf{Mid training} ($\mu_n$ decreasing, $V_i$ decreasing):
    \begin{itemize}
      \item L1 relaxes → policy support grows ($k$ increases)
      \item But evidence is improving (more data) → $V_i$ decreases
      \item Net effect: $d_{\text{eff}}$ grows but $V_i^{\text{eff}}$ stays controlled
      \item Agent explores new actions with calibrated confidence
    \end{itemize}

  \item \textbf{Late training} ($\mu_n \approx 0$, $V_i$ small):
    \begin{itemize}
      \item No sparsity constraint → full action support available
      \item High evidence → agent knows its policy well
      \item Converges to true (possibly dense) Nash equilibrium
    \end{itemize}
\end{enumerate}

\begin{remark}[Comparison with entropy regularization]
Standard entropy regularization $-\beta H(\pi)$ encourages \emph{uniform}
exploration (spread mass). L1 encourages \emph{focused} exploration
(concentrate mass). In high dimensions, L1 is strictly better:
\begin{itemize}
  \item Entropy: $d_{\text{eff}} \approx d$ (all dimensions active) → no blessing
  \item L1: $d_{\text{eff}} \approx k$ → full blessing of dimensionality
\end{itemize}
This is why entropy regularization helps with exploration but
\emph{does not} address the variance problem, while L1 addresses both.
\end{remark}


\section{Convergence Analysis}

\subsection{Main Result}

\begin{theorem}[Convergence of Sparse-EW-PG]\label{thm:main}
Let $\pi^*$ be a $k$-sparse stable Nash policy in a stochastic game
with $|\mathcal{A}| = d$. Under the SOS condition with parameter $\mu$,
and with:
\begin{itemize}
  \item $\gamma_n = \gamma/(n+m)^p$, $p \in (1/2, 1]$
  \item $\mu_n = \mu_0 / (1 + n/n_0)$ (L1 annealing)
  \item Evidence weights computed with effective dimension
\end{itemize}
the Sparse-EW-PG iterates satisfy:
\begin{equation}\label{eq:main-result}
  \E[\norm{\pi_n - \pi^*}^2] \;=\; O\!\left(
    \frac{\HM(V^{\text{eff}})}{\AM(V^{\text{eff}})} \cdot
    \frac{C_k}{n^q}
  \right)
\end{equation}
where $C_k = O(k \log(d/k))$ is the \emph{sparse constant},
sub-linear in~$d$.
\end{theorem}

\begin{proof}[Proof sketch]
Three steps:
\begin{enumerate}
  \item The L1 proximal step ensures $\pi_n$ stays approximately
    $k$-sparse for all~$n$ (up to the annealing schedule). Standard
    proximal gradient convergence gives $\norm{\pi_n - \pi^*_k}_1 \to 0$
    where $\pi^*_k$ is the best $k$-sparse approximation to $\pi^*$.

  \item With $k$-sparse iterates, the effective dimension satisfies
    $d_{\text{eff},n} \leq 2k$ (the gradient covariance has at most
    $2k$ significant eigenvalues --- $k$ from the policy support
    and $k$ from the exploration directions).

  \item Apply the EW-PG convergence theorem with effective variance:
    $V^{\text{eff}} = (d_{\text{eff}}/d) \cdot V$.
    The constant $C_k$ arises from \cref{thm:cs}: the sparse recovery
    bound replaces the ambient-dimension bound.
\end{enumerate}
\end{proof}

\subsection{Comparison of Sample Complexities}

\begin{center}
\begin{tabular}{lccc}
\hline
\textbf{Method} & \textbf{Samples to $\epsilon$-Nash} & \textbf{Scaling in $d$} & \textbf{Scaling in $k$} \\
\hline
Standard PG             & $O(d^2 / \epsilon^2)$                & $d^2$ (quadratic)     & --- \\
PG + entropy reg.       & $O(d / \epsilon^2)$                  & $d$ (linear)          & --- \\
PG + L1 reg.            & $O(k \log(d/k) / \epsilon^2)$        & $\log d$ (logarithmic) & $k$ (linear) \\
EW-PG + L1 (ours)       & $O(\frac{\HM}{\AM} k \log(d/k) / \epsilon^2)$ & $\log d$ & $k \cdot \frac{\HM}{\AM}$ \\
\hline
\end{tabular}
\end{center}

The combined method achieves \emph{logarithmic} scaling in the action space
dimension and an additional $\HM/\AM \leq 1$ improvement from evidence weighting.


\section{The Communication Connection}

\subsection{Sparse Policies Are Communicable}

The cooperative PG (companion paper) shows that communication quality
is bounded by self-knowledge, which is bounded by evidence. The
blessing of dimensionality adds a crucial refinement:
\emph{sparse policies are inherently more communicable}.

To communicate a $k$-sparse policy over a $d$-action space:
\begin{itemize}
  \item Naive: transmit $d$ probabilities → rate $O(d)$
  \item Sparse: transmit $k$ (action, probability) pairs → rate $O(k \log d)$
\end{itemize}

The $\log d$ factor (identifying which actions) replaces the $d$ factor
(transmitting all probabilities). This means:

\begin{corollary}[Sparsity enables coalition formation]
In a game with $|\mathcal{A}| = d$ and $k$-sparse optimal policies,
the communication bandwidth required for effective coalition coordination
scales as $O(k \log d)$, not $O(d)$. The self-knowledge bottleneck
(\S2 of Cooperative PG) is relaxed by a factor of $d / (k \log d)$.
\end{corollary}

\subsection{Vibing Is Dense; Articulation Is Sparse}

The ``vibing'' spectrum from the cooperative paper acquires a precise meaning:
\begin{itemize}
  \item \textbf{Vibing agent}: policy is effectively dense (spreading mass
    over many actions through imprecise habits). $d_{\text{eff}} \approx d$.
    Poor self-knowledge, poor communication.
  \item \textbf{Articulate agent}: policy is sparse (concentrating on the
    actions that matter). $d_{\text{eff}} \approx k \ll d$.
    Good self-knowledge, good communication.
\end{itemize}

L1 regularization drives the agent from vibing toward articulation.
The blessing of dimensionality ensures this transition is
\emph{computationally cheap} in the action space dimension.


\section{Connection to the $\Omega$-Framework}

\subsection{The Embedding Space Interpretation}

In the $\Omega$-framework (theos \S I), each agent's observation space
$\mathcal{O}$ is embedded in a high-dimensional space where the blessing
of dimensionality operates. The manuscript (Chapter~9) describes this
as ``Maximus the Confessor's logoi --- every entity has a position in
semantic space,'' and notes that ``in high dimensions, structure becomes
MORE visible, not less.''

The policy gradient result makes this precise:
\begin{itemize}
  \item The policy $\pi \in \R^d$ lives in a high-dimensional space
  \item The \emph{effective} policy lives on a $k$-dimensional submanifold
  \item The gradient concentrates on this submanifold (blessing)
  \item The evidence weight measures information content on the submanifold
  \item Sparse recovery lets us find the submanifold from few samples
\end{itemize}

This is the ``blessing of dimensionality'' from Chapter~9 of the
theological manuscript, now formalized as a convergence guarantee.

\subsection{Pearl's Hierarchy and Sparse Representations}

At Pearl Level 1 (association), the agent observes correlations across
all $d$ actions --- a dense, high-variance representation. At Pearl
Level 3 (counterfactual), the agent has identified the $k$ \emph{causal}
actions --- a sparse, low-variance representation.

The progression up Pearl's hierarchy IS the sparsification process:
\begin{center}
Level 1: dense policy ($d_{\text{eff}} \approx d$, high $V$, ``vibing'') \\
$\downarrow$ \quad {\small L1 regularization + evidence accumulation} \\
Level 3: sparse policy ($d_{\text{eff}} \approx k$, low $V^{\text{eff}}$, ``articulate'')
\end{center}

The L1 coefficient $\mu_n$ controls the speed of this ascent.


\section{Numerical Illustration}

\subsection{Expected Behavior}

For a game with $|\mathcal{A}| = 1000$ and $k = 5$ optimal actions:

\begin{center}
\begin{tabular}{lccc}
\hline
\textbf{Phase} & $d_{\text{eff}}$ & $V^{\text{eff}} / V$ & \textbf{Effective samples} \\
\hline
Early ($n < 100$)     & $\sim$3--5   & $\sim$0.5\%   & $200\times$ more informative \\
Mid ($100 < n < 1000$) & $\sim$5--15  & $\sim$1.5\%  & $65\times$ more informative \\
Late ($n > 1000$)     & $\sim$5--30  & $\sim$3\%     & $33\times$ more informative \\
\hline
\end{tabular}
\end{center}

The ``effective samples'' column shows how much more informative each
sample is compared to the ambient-dimension estimate. Even in the late
phase, the effective dimension is $30 \ll 1000$, giving a $33\times$
efficiency gain from the blessing of dimensionality alone.

\appendix
\section{Background: Key Results from High-Dimensional Statistics}

For reference, the main results we invoke:

\begin{theorem}[Matrix Bernstein inequality, Tropp 2012]
Let $X_1, \ldots, X_M$ be independent random matrices with $\E[X_m] = 0$
and $\norm{X_m} \leq L$. Then:
\[
  \Prob\!\left(\left\|\sum_m X_m\right\| \geq t\right)
  \leq (d_1 + d_2) \exp\!\left(\frac{-t^2/2}{\sigma^2 + Lt/3}\right)
\]
where $\sigma^2 = \max(\norm{\sum \E[X_m X_m^\top]}, \norm{\sum \E[X_m^\top X_m]})$.
\end{theorem}

\begin{theorem}[Cand\`es--Tao RIP, 2005]
Let $\Phi \in \R^{M \times d}$ satisfy the restricted isometry property
of order $2k$: for all $2k$-sparse $x$,
$(1-\delta_{2k})\norm{x}^2 \leq \norm{\Phi x}^2 \leq (1+\delta_{2k})\norm{x}^2$.
If $\delta_{2k} < \sqrt{2} - 1$, then $\ell_1$ minimization recovers
any $k$-sparse signal exactly from $M = O(k \log(d/k))$ measurements.
\end{theorem}

\begin{theorem}[Johnson--Lindenstrauss, 1984]
For any $\epsilon > 0$ and $n$ points in $\R^d$, there exists a linear
map $\Phi : \R^d \to \R^m$ with $m = O(\log n / \epsilon^2)$ preserving
all pairwise distances within factor $(1 \pm \epsilon)$.
\end{theorem}


\end{document}
